Photosynthesis is the most vital process on Earth, sustaining life by means of food, oxygen, and other resources. This importance, however, brings forth its own challenges. It is a complex process, with many interacting components and feedback loops. This complexity has yet to deter humans from trying to understand it. Over the past centuries, photosynthesis has been a main focus of biological research. However, the scope of scientific work has since plateaued. To further explore the possibilities, more intricacies and details of the process need to be understood, and one of the most prevalent tools for doing so is mathematical modeling. These models are not new inventions, but they have grown in popularity and complexity over the years. However, the current state of the art is not without its flaws. Due to the complexity of photosynthesis, models narrow their scope to make them more manageable. This reduction, however, leads to many possible models. Over the years, the number of publications that have mentioned photosynthesis models has been ever-increasing. This increase in models, consequently, also brings a complication: the risk of being overwhelmed by the choice. There have been attempts to solve this problem, but they have only been partially successful. These solutions aim to solve problems across as many models as possible, but that fact undermines their capabilities. The models are so vastly different that many shortcuts are needed to incorporate them all into a single framework. Conversely, embracing model diversity and focusing on specific model types reduces the scope of the solution. By doing so, the light can shine on the specific type of models. More care can be taken to ensure these models are kept to a standard that follows modern guidelines like \glsentryshort{fair} and \glsentryshort{cure}. On top of that, it allows a more direct comparison, so direct that standardized demonstrations can be used to showcase the models' capabilities. Moreover, providing a web platform to share these models with the community helps ensure they are more easily found and used. These three aspects are the core of \greensloth{}, which comprises an entire ecosystem of tools for creating, demonstrating, and sharing \glsentryshort{ode} models of photosynthesis. The website (\url{https://greensloth.rwth-aachen.de/}) is the main platform for sharing the models, and the tools for creating and demonstrating them are also available as open-source GitHub repositories (\url{https://github.com/ElouenCorvest/GreenSloth/tree/main}). To highlight the capabilities of \greensloth{}, five models of photosynthesis are used to showcase the ecosystem's strengths and weaknesses. These models vary in the scope of photosynthesis and in complexity, providing a good starting point as a proof of concept.